\documentclass[11pt]{article}
\usepackage[margin=1in]{geometry}
\usepackage[T1]{fontenc}
\usepackage[utf8]{inputenc}
\usepackage{amsmath,amssymb,amsthm}
\usepackage{enumitem}

\title{ICPC World Finals 2020\\G. Opportunity Cost}
\author{}
\date{}

\begin{document}
\maketitle

\section*{Problem Summary}

Each phone has three scores $(x,y,z)$.
If we buy phone $(x,y,z)$, then another phone $(x_i,y_i,z_i)$ can dominate it by
\[
    \max(x_i-x,0)+\max(y_i-y,0)+\max(z_i-z,0).
\]
The opportunity cost of our phone is the maximum of this quantity over all phones.

We must find the phone with minimum opportunity cost, breaking ties by smallest index.

\section*{Key Observations}

\begin{itemize}[leftmargin=*]
    \item For a fixed phone $(x,y,z)$ and a fixed competitor $(x_i,y_i,z_i)$,
    \[
        \max(x_i-x,0)+\max(y_i-y,0)+\max(z_i-z,0)
    \]
    is the sum of the positive coordinate differences.
    \item That is exactly
    \[
        \max_{S \subseteq \{x,y,z\}}
        \left( \sum_{d \in S} d_i - \sum_{d \in S} d \right),
    \]
    because for each coordinate we either include its difference (if positive) or exclude it (if negative).
    \item Therefore the opportunity cost of phone $(x,y,z)$ is
    \[
        \max_{\emptyset \ne S \subseteq \{x,y,z\}}
        \left( M_S - \sum_{d \in S} d \right),
    \]
    where
    \[
        M_S = \max_i \sum_{d \in S} d_i.
    \]
    \item There are only $2^3-1=7$ nonempty subsets, so we can precompute all values $M_S$ in one pass and
    then evaluate every phone in constant time.
\end{itemize}

\section*{Algorithm}

\begin{enumerate}[leftmargin=*]
    \item Number the three dimensions by bits $0,1,2$.
    For every nonempty mask \texttt{mask} from $1$ to $7$, let
    \[
        M_{\texttt{mask}} = \max_i \sum_{\texttt{bit} \in \texttt{mask}} \text{value}_{i,\texttt{bit}}.
    \]
    \item Read all phones and update these $7$ maxima.
    \item For each phone $p$, compute
    \[
        \max_{\texttt{mask}=1}^{7}
        \left( M_{\texttt{mask}} - \sum_{\texttt{bit} \in \texttt{mask}} p_{\texttt{bit}} \right).
    \]
    This is exactly its opportunity cost.
    \item Keep the minimum pair
    \[
        (\text{cost}, \text{index})
    \]
    lexicographically, so ties are broken by the smallest index automatically.
\end{enumerate}

\section*{Correctness Proof}

We prove that the algorithm returns the correct answer.

\paragraph{Lemma 1.}
For any real numbers $a,b,c$,
\[
    \max(a,0)+\max(b,0)+\max(c,0)
    =
    \max_{S \subseteq \{1,2,3\}} \sum_{j \in S} v_j,
\]
where $(v_1,v_2,v_3)=(a,b,c)$.

\paragraph{Proof.}
For each coordinate $v_j$, adding it helps exactly when $v_j > 0$ and hurts when $v_j < 0$.
So the maximum subset sum is obtained by taking precisely the positive coordinates, and its value is
$\max(a,0)+\max(b,0)+\max(c,0)$. \qed

\paragraph{Lemma 2.}
For a fixed phone $p=(x,y,z)$, its opportunity cost equals
\[
    \max_{\emptyset \ne S \subseteq \{x,y,z\}}
    \left( M_S - \sum_{d \in S} d \right).
\]

\paragraph{Proof.}
For any competitor phone $i$, apply Lemma 1 to
\[
    a = x_i-x,\qquad b = y_i-y,\qquad c = z_i-z.
\]
Then
\[
    \max(x_i-x,0)+\max(y_i-y,0)+\max(z_i-z,0)
    =
    \max_S \left( \sum_{d \in S} d_i - \sum_{d \in S} d \right).
\]
Now maximize over all competitors $i$:
\[
    \max_i \max_S (\cdots)
    =
    \max_S \max_i (\cdots)
    =
    \max_S \left( M_S - \sum_{d \in S} d \right).
\]
The empty subset contributes $0$, so restricting to nonempty subsets does not change the maximum. \qed

\paragraph{Lemma 3.}
For every phone $p$, the algorithm computes its exact opportunity cost.

\paragraph{Proof.}
The preprocessing step computes every value $M_S$ exactly by definition.
Then for phone $p$ the algorithm evaluates the expression from Lemma 2 over all seven nonempty subsets.
Hence the maximum it obtains is exactly the opportunity cost of $p$. \qed

\paragraph{Theorem.}
The algorithm outputs the correct answer for every valid input.

\paragraph{Proof.}
By Lemma 3, the algorithm computes the exact opportunity cost of every phone.
It then chooses the minimum pair $(\text{cost}, \text{index})$, so it returns the smallest possible
opportunity cost and, among all phones achieving it, the smallest index.
Therefore the output is correct. \qed

\section*{Complexity Analysis}

There are only $7$ nonempty subsets.
So both the preprocessing pass and the evaluation pass take $O(7n)=O(n)$ time.

The stored phone list uses $O(n)$ memory.

\section*{Implementation Notes}

\begin{itemize}[leftmargin=*]
    \item Using bitmasks $1 \dots 7$ is a convenient way to enumerate the seven nonempty subsets.
    \item All sums fit safely in 64-bit integers.
    \item The lexicographic minimum of \texttt{(cost, index)} handles the required tie-break automatically.
\end{itemize}

\end{document}
